\section*{Problems}

\subsection*{Problem statements}

% Remember
\begin{problema}
	\label{prob:de3bea1}
	Classify each of the following variables according to its scale of measurement (\emph{Nominal, Ordinal, Interval, or Ratio}):
	\begin{enumerate}
		\item The postal code of a customer's residence on an e-commerce platform.
		\item A user's satisfaction level with a mobile app (Very Dissatisfied, Dissatisfied, Neutral, Satisfied, Very Satisfied).
		\item The temperature in degrees Celsius (°C) recorded by the servers in a data center.
		\item The monthly salary in Mexican pesos of a data science engineer.
	\end{enumerate}
\end{problema}

% Understand
\begin{problema}
	\label{prob:7b7dc88}
	Determine whether the metric described in each statement is a \textbf{Parameter} ($\mu, \sigma^2, p$) or a \textbf{Statistic} ($\bar{X}, S^2, \hat{p}$), and explain the criterion that distinguishes the two cases:
	\begin{enumerate}
		\item After surveying all 32 federal entities, INEGI determined that the national average number of people per household is $3.6$.
		\item A quality engineer selects a random batch of 80 processors and computes a clock-frequency variance of $0.015 \text{ GHz}^2$.
		\item In a test subset of 5,000 banking transactions, the detected fraud rate was $0.84\%$.
	\end{enumerate}
\end{problema}

% Apply
\begin{problema}
	\label{prob:4bc23ad}
	A technology consulting corporation has 4,000 employees divided into four divisions: Development ($2,000$), Infrastructure ($1,000$), Sales ($600$), and Human Resources ($400$). A random sample of $n = 200$ employees must be selected to evaluate the work climate. Design a \emph{proportional stratified random sampling} scheme and determine exactly how many employees should be interviewed in each of the four divisions.
\end{problema}

% Analyze
\begin{problema}
	\label{prob:a5bddfb}
	In 1936, \emph{Literary Digest} conducted a massive survey by mailing 10 million questionnaires to predict the United States presidential election, using telephone directories and automobile-owner records as the sampling frame. Despite obtaining a gigantic sample of $n \approx 2.4$ million responses, they predicted a decisive victory for Alf Landon, when Roosevelt actually won in a landslide. Show conceptually and algebraically, using the decomposition $\text{MSE}(\bar{X}_n) = \sigma^2/n + B^2$, how a systematic bias $B = \mathbb{E}[\bar{X}] - \mu \neq 0$ makes increasing $n \to \infty$ fail to reduce total error; instead, the estimate converges with certainty to an incorrect value.
\end{problema}

% Evaluate
\begin{problema}
	\label{prob:feeec4c}
	To evaluate the success of a new web-banking platform, a user-experience department sends a mass email to 100,000 customers inviting them to voluntarily answer an online usability questionnaire. A total of 1,200 customers respond, and $88\%$ of them state that the platform is extremely confusing. Evaluate which methodological bias predominates in this data collection, why it distorts the true proportion among all customers ($p$), and propose a redesign of the sampling protocol that would yield an unbiased estimate without increasing the collection budget.
\end{problema}

% Create
\begin{problema}
	\label{prob:701b535}
	Design an original example (different from the \emph{Literary Digest}, streaming, or web-banking cases already discussed) of a data science study in which the collection mechanism introduces systematic coverage or self-selection bias. Explicitly identify the target population, the sample actually obtained, the type of bias present, and a redesigned sampling protocol that eliminates or mitigates it.
\end{problema}

\subsection*{Detailed solutions}

% Remember
\begin{solproblema}[prob:de3bea1]
	The classification by measurement level or scale is justified as follows:
	\begin{enumerate}
		\item \textbf{Postal code $\to$ Nominal:} The numbers act only as labels or categories for geographic location. There is no natural order and no mathematical meaning in adding or subtracting postal codes.
		\item \textbf{User satisfaction $\to$ Ordinal:} There is an intrinsic ranking (Very Dissatisfied $<$ Dissatisfied $<$ Neutral $<$ Satisfied $<$ Very Satisfied), but the conceptual distance between consecutive categories is not necessarily uniform or quantifiable.
		\item \textbf{Temperature in °C $\to$ Interval:} A numerical difference in degrees has a constant physical meaning, but zero is arbitrary ($0^\circ\text{C}$ represents the freezing point of water, not the absolute absence of thermal energy). Therefore, we cannot say that $40^\circ\text{C}$ is "twice as hot" as $20^\circ\text{C}$.
		\item \textbf{Monthly salary $\to$ Ratio:} It has uniform intervals and a real absolute zero ($0$ pesos means a complete absence of salary income). Ratios are valid: a salary of $\$60,000$ is exactly twice one of $\$30,000$.
	\end{enumerate}
\end{solproblema}

% Understand
\begin{solproblema}[prob:7b7dc88]
	The distinction between parameter and statistic depends on whether the metric describes the entire universe being studied or only a portion of it:
	\begin{enumerate}
		\item \textbf{Average of $3.6$ people per household $\to$ Parameter ($\mu$):} It comes from an exhaustive national census that covered all 32 federal entities.
		\item \textbf{Variance of $0.015 \text{ GHz}^2$ $\to$ Statistic ($S^2$):} It is a measure computed only from the random sample of 80 inspected processors.
		\item \textbf{Fraud rate of $0.84\%$ $\to$ Statistic ($\hat{p}$):} It is a proportion estimated from the sample subset of 5,000 transactions extracted for testing.
	\end{enumerate}
	The decisive criterion is not the size of the data set, but whether it came from an exhaustive census (parameter) or from a random portion of the universe (statistic).
\end{solproblema}

% Apply
\begin{solproblema}[prob:4bc23ad]
	In proportional stratified random sampling, the sample size for each stratum $k$ is proportional to its weight in the payroll: $n_k = n \cdot \left(\frac{N_k}{N}\right) = 200 \cdot \left(\frac{N_k}{4,000}\right)$. Substituting:
	\begin{itemize}
		\item \textbf{Development:} $n_1 = 200 \cdot \left(\frac{2,000}{4,000}\right) = 100$ employees.
		\item \textbf{Infrastructure:} $n_2 = 200 \cdot \left(\frac{1,000}{4,000}\right) = 50$ employees.
		\item \textbf{Sales:} $n_3 = 200 \cdot \left(\frac{600}{4,000}\right) = 30$ employees.
		\item \textbf{Human Resources:} $n_4 = 200 \cdot \left(\frac{400}{4,000}\right) = 20$ employees.
	\end{itemize}
	The total verifies that $100 + 50 + 30 + 20 = 200$ interviews, with structural parity guaranteed relative to the true weight of each division.
\end{solproblema}

% Analyze
\begin{solproblema}[prob:a5bddfb]
	Let $\bar{X}_n$ be the sample estimator based on a sample of size $n$. The mean squared error (MSE) of an estimator decomposes total imprecision into sampling variance and squared bias:
	\begin{align}
		\text{MSE}(\bar{X}_n) = \mathbb{E}\left[(\bar{X}_n - \mu)^2\right] = \text{Var}(\bar{X}_n) + (\text{Bias})^2 = \dfrac{\sigma^2}{n} + B^2,
	\end{align}
	where $B = \mathbb{E}[\bar{X}_n] - \mu$ represents the methodological bias induced by the selection frame (in \emph{Literary Digest}, surveying only wealthier people with telephones and automobiles during the Great Depression).

	As the sample size increases indefinitely ($n \to \infty$), the sampling-variability component tends to zero ($\lim_{n \to \infty} \sigma^2/n = 0$), by the Law of Large Numbers. However, the systematic-bias component remains exactly the same:
	\begin{align}
		\lim_{n \to \infty} \text{MSE}(\bar{X}_n) = 0 + B^2 = B^2 > 0.
	\end{align}
	This shows that a gigantic sample of millions of biased records ($n \approx 2.4$ million) merely reduces variance around an incorrect value, creating a false illusion of precision for a wrong estimate. By contrast, a smaller sample such as Gallup's ($n = 50,000$), where $B \approx 0$, can converge without bias toward the true parameter $\mu$.
\end{solproblema}

% Evaluate
\begin{solproblema}[prob:feeec4c]
	The dominant issue is \textbf{self-selection or voluntary-response bias}. Users who experienced severe technical frustration with the platform have much stronger psychological motivation to open the survey and complain than users who were satisfied or indifferent and completed their operations without problems. Therefore, the $88\%$ figure is a massive overestimate of the true population level of confusion ($p$).

	\textbf{Redesign with no extra cost:} instead of sending a passive email to 100,000 people, select a much smaller purely random sample (for example, $n = 500$ customers) and apply active intermittent follow-up in the same web session (a controlled random pop-up after completing a transaction) or scheduled direct contact, ensuring a high response rate that is not biased by personal frustration.
\end{solproblema}

% Create
\begin{solproblema}[prob:701b535]
	\textbf{Example:} a university wants to estimate the average weekly hours that its graduates dedicate to continuing education, and it sends the survey only through its official LinkedIn group. \textbf{Target population:} all graduates of the university. \textbf{Obtained sample:} only graduates who keep an active LinkedIn account and follow the institutional page---a subset biased toward those who are more professionally active and connected. \textbf{Type of bias:} coverage bias (the sampling frame---LinkedIn users---systematically excludes graduates without a presence on that network, typically those who changed careers or became disconnected from the formal labor market), worsened by possible self-selection bias among those who do see the post. \textbf{Redesign:} use the complete institutional graduate directory (the email address registered at graduation) as the sampling frame, and draw a simple random sample from it instead of depending on voluntary membership in a social network.
\end{solproblema}
